% ============================================================
%  6.1 Fundamentos del dominio   (capítulo 6, Marco teórico)
%  Se incluye desde secciones/07_marco_teorico.tex con \input.
% ============================================================
\section{Fundamentos del dominio}
\label{sec:fundamentos-dominio}

\subsection{Transporte público informal (paratránsito) en ciudades latinoamericanas y del Sur Global}
\label{subsec:paratransito}

Se denomina transporte informal o paratránsito al conjunto de servicios de transporte de
pasajeros de tipo paratránsito que se prestan sin sanción oficial plena, generalmente con
vehículos de pequeña y mediana capacidad operados por propietarios individuales agrupados
en sindicatos, cooperativas o asociaciones \parencite{cervero2007}. En su revisión de
experiencias internacionales, \textcite{cervero2007} señalan que estos sistemas ofrecen
beneficios reales, como la movilidad a demanda para la población dependiente del transporte
público, el empleo para trabajadores de baja calificación y la cobertura de áreas carentes
de oferta formal, pero también generan costos, como congestión, contaminación atmosférica y
acústica, y accidentes de tránsito. El paratránsito no constituye, por tanto, una anomalía
marginal, sino la forma dominante de movilidad colectiva en numerosas ciudades de América
Latina, África y Asia, lo que obliga a analizarlo con herramientas propias en lugar de
asumir las categorías del transporte formal.

Uno de los rasgos que más condiciona el análisis de estos sistemas es su problema de
información. Las rutas suelen ser conocidas por los usuarios habituales, pero no están
documentadas en fuentes oficiales, y los vehículos pueden detenerse en cualquier punto del
recorrido. La propia Trufi Association describe este tipo de servicio como minibuses que
llevan a los pasajeros a su destino por rutas conocidas pero no documentadas, sin paradas
fijas, porque los pasajeros pueden subir y bajar en cualquier punto mientras haya espacio
\parencite{trufi2019}. \textcite{williams2015} constatan, para el caso de Nairobi, que en
muchas ciudades del Sur Global prácticamente no existe información digital sobre rutas,
paradas, frecuencias ni horarios de los minibuses semiformales. En Cochabamba, en
particular, no existía un archivo GTFS de transporte gestionado por la administración
pública, por lo que la información de rutas debió generarse de forma colaborativa y
publicarse como datos abiertos \parencite{osor-trufi}.

El caso de Cochabamba ilustra con claridad esta situación. \textcite{cabrera2022} sostienen
que en las ciudades bolivianas la mayor parte de los servicios de movilidad de pasajeros
pueden clasificarse como paratránsito, y caracterizan la región metropolitana de Cochabamba,
con más de 1,2~millones de habitantes en áreas urbanas de varios municipios conurbados, como
un sistema operado por centenas de organizaciones pequeñas y medianas, en el que, según una
nota del mismo artículo, cerca de 840\,000 personas realizan 2~millones de viajes diarios,
el 75~\% de ellos en el eje Sacaba-Cochabamba-Quillacollo. Según los datos que recogen estos
autores, el 55~\% de los desplazamientos se realiza en transporte público o paratránsito, el
26~\% a pie o en bicicleta, el 16~\% en vehículo privado y el 4~\% en taxi; asimismo, un
inventario realizado entre 2017 y 2020 por estudiantes de urbanismo de la Universidad Mayor
de San Simón y la Universidad Privada Boliviana identificó 132 líneas con 648 rutas
diferentes (45 de coaster, 63 de micros, 193 de taxitrufis y 347 de trufis). Los mismos
autores describen cuatro tipos de vehículo (micros, coaster, trufis y taxitrufis) y tres
formas de agencia sobre las rutas, que son la creación, la extensión y la subdivisión, a
través de las cuales la oferta acompaña, e incluso induce, la expansión urbana informal.

La implicación para el proyecto es doble. Por un lado, la ausencia de datos oficiales de
rutas convierte a la cartografía colaborativa en la única fuente sistemática de oferta, con
la consecuencia de que la cobertura observable depende del esfuerzo de mapeo y no solo de la
existencia física del servicio. Por otro lado, el carácter dinámico de las rutas, que se
crean, extienden y subdividen, implica que cualquier representación de la oferta es una
instantánea con vigencia limitada, lo que justifica un enfoque orientado a priorizar dónde
revisar la cartografía antes que a certificar su completitud.

\subsection{Cartografía colaborativa del transporte informal: OpenStreetMap, GTFS en sistemas semiformales y Trufi}
\label{subsec:cartografia-colaborativa}

OpenStreetMap (OSM) es un proyecto de cartografía abierta en el que voluntarios crean y
mantienen una base de datos geográfica mundial bajo licencia abierta \parencite{haklay2008}.
Se inscribe en el fenómeno que \textcite{goodchild2007} denominó información geográfica
voluntaria, según el cual ciudadanos sin formación cartográfica formal actúan como sensores
y productores de datos espaciales. En OSM, las líneas de transporte público se representan
mediante relaciones que agrupan las vías recorridas y, opcionalmente, las paradas, lo que
permite describir rutas de paratránsito aun cuando carezcan de paradas oficiales. La calidad
y la completitud de esta información son, sin embargo, heterogéneas en el espacio, porque
dependen de la distribución geográfica del esfuerzo voluntario.

El antecedente de referencia en la digitalización del paratránsito es el proyecto Digital
Matatus de Nairobi. \textcite{williams2015} muestran que la tecnología de telefonía celular,
ubicua en la mayoría de los países, puede emplearse para generar un conjunto de datos del
sistema semiformal de matatus en un estándar abierto, el GTFS (\textit{General Transit Feed
Specification}), concebido originalmente para el transporte formal. Los autores concluyen
que el GTFS puede adaptarse a sistemas semiformales, que existe demanda de datos
estandarizados sobre estos sistemas en los países en desarrollo y que la publicación de los
datos en un estándar abierto estimula el desarrollo de tecnologías como las aplicaciones de
ruteo. \textcite{klopp2015} complementan este trabajo con las lecciones del mismo proyecto
para el diseño de herramientas de orientación y planificación de viajes basadas en teléfonos
móviles.

En Cochabamba, esta función la cumple la Trufi Association, una organización sin fines de
lucro que desarrolla aplicaciones de movilidad para el transporte informal, que lanzó su
aplicación en Cochabamba en febrero de 2019 y que se fundó oficialmente en Hamburgo
(Alemania) en abril de ese mismo año \parencite{trufi2019}. Según la propia asociación,
desde el inicio del proyecto en abril de 2018 un equipo local recopiló más de 200 rutas de
los trufis de Cochabamba, primero en un sistema propio y luego como datos públicos en
OpenStreetMap, y en agosto de 2019 el código de la aplicación se liberó como software de
código abierto \parencite{trufi2019}. La cadena de procesamiento documentada en los
repositorios de la asociación extrae las relaciones de transporte de OSM, las convierte en
GeoJSON, genera a partir de ellas un archivo GTFS y lo carga en OpenTripPlanner para
calcular itinerarios. Versiones posteriores de la aplicación incorporaron nuevas rutas
añadidas mediante OSM y una función para que los usuarios reporten rutas faltantes o
incorrectas \parencite{trufi-v3}; la publicación de la versión~3 anuncia más de cien rutas
nuevas, mientras que otra publicación de trufi.app informa de 25 rutas nuevas con 104
conexiones nuevas, para un total de 120 rutas y 339 conexiones en Cochabamba y sus
alrededores, lo que sugiere que parte de esas cifras corresponde a conexiones y no a rutas
\parencite{trufiapp-25rutas}. La asociación informó haber superado las 20\,000 instalaciones
en Android, cerca de 30\,000 sumando las de iOS \parencite{trufi-20000}, y una publicación
posterior eleva la cifra a 100\,000 instalaciones en Android y cerca de 25\,000 en iOS, salto
que atribuye a un video de TikTok de Comteco, una cooperativa de telecomunicaciones de
Cochabamba \parencite{trufi-100000}. En su versión~5, la aplicación ofrece rutas mantenidas
por voluntarios que mapean en OpenStreetMap, con ruteo sin conexión basado en GTFS
\parencite{trufi-v5}; la ficha de la aplicación en App Store precisa 626 rutas, incluidas las
líneas Verde y Roja y la del teleférico \parencite{trufi-appstore}. Esa cifra corresponde a
recorridos y no a líneas: el feed empleado en este trabajo describe 141 líneas y 626
recorridos, porque una línea puede tener varios recorridos, por ejemplo ida y vuelta
(sección~\ref{sec:fuentes-datos}).

De lo anterior se desprende que la utilidad de ``Trufi App'' en cada lugar depende directamente
de que las rutas que lo sirven hayan sido mapeadas. Una zona sin rutas cargadas no puede
producir respuestas útiles para el usuario, lo que puede desalentar su uso y reducir las
consultas; por ello, el patrón espacial de consultas contiene información sobre el esfuerzo
cartográfico, y el problema de priorizar dónde revisar o mapear rutas es, en sentido
estricto, un problema de asignación de un recurso voluntario escaso.

\subsection{El formato GTFS: estructura, trazado frente a paradas y limitaciones en sistemas informales}
\label{subsec:gtfs}

La especificación GTFS Schedule, hoy mantenida por la organización MobilityData, define un
conjunto de archivos de texto delimitados por comas que describen la oferta programada de
un sistema de transporte público \parencite{mobilitydata-gtfs}. Entre sus archivos centrales
se encuentran \var{agency.txt} (operadores), \var{routes.txt} (líneas o rutas, con su tipo
de vehículo), \var{trips.txt} (viajes concretos de cada ruta, asociados a un calendario de
servicio y, opcionalmente, a un trazado), \var{stops.txt} (paradas o puntos de acceso con sus
coordenadas), \var{stop_times.txt} (secuencia y horarios de paso de cada viaje por las
paradas) y \var{shapes.txt}, archivo opcional que describe el trazado geográfico que sigue el
vehículo mediante una secuencia ordenada de puntos (\var{shape_id}, \var{shape_pt_lat},
\var{shape_pt_lon}, \var{shape_pt_sequence}). El estándar nació de la colaboración entre la
agencia de transporte de Portland (TriMet) y Google y se convirtió en el formato de facto
para el intercambio de datos de transporte público \parencite{mchugh2013}.

Para el presente proyecto resulta fundamental distinguir entre trazado y paradas. El trazado
(\var{shapes.txt}) representa la geometría lineal por la que circula el vehículo, mientras
que las paradas (\var{stops.txt}) representan los puntos discretos en los que el estándar
supone que se produce el ascenso y el descenso. En un sistema formal ambas representaciones
coinciden en lo sustancial, porque el acceso ocurre solo en paradas. En el paratránsito
cochabambino, en cambio, el ascenso y el descenso pueden producirse a lo largo de casi todo
el recorrido, de modo que las paradas incluidas en un GTFS generado a partir de OSM son, en
gran medida, convencionales: puntos que el generador del feed ubica sobre el recorrido para
satisfacer la estructura del estándar, y no paradas reconocidas oficialmente. Medir el
acceso como distancia a esas paradas introduciría, por tanto, un artefacto del proceso de
generación del feed.

El GTFS presenta además limitaciones conocidas cuando se aplica a sistemas informales. La
primera es la vigencia: el feed es una instantánea de las rutas mapeadas en la fecha de su
generación, y las rutas del paratránsito se crean, extienden y subdividen con frecuencia
\parencite{cabrera2022}. La segunda es la naturaleza aproximada de los horarios y
frecuencias, que en ausencia de tablas oficiales se derivan de supuestos de velocidad y de
intervalos nominales. La tercera, y más relevante, es que el feed representa la oferta
mapeada y no la oferta real, porque una ruta existente pero no mapeada simplemente no
aparece. De ello se deriva que el proyecto emplea el trazado (\var{shapes.txt}) como
representación más fiel de la oferta en un sistema sin paradas fijas, y que interpreta
cualquier medida de cobertura como cobertura de la oferta conocida por ``Trufi App'', no como
accesibilidad efectiva al transporte.

\subsection{Cobertura y accesibilidad al transporte público: ODS 11.2.1 y umbrales de caminata}
\label{subsec:ods-11-2-1}

La Agenda 2030 incorpora en su meta 11.2 el acceso a sistemas de transporte seguros,
asequibles, accesibles y sostenibles, y la monitorea mediante el indicador 11.2.1,
\enquote{proporción de la población que tiene acceso conveniente al transporte público,
desglosada por sexo, edad y personas con discapacidad}, cuyo organismo custodio es el
Programa de las Naciones Unidas para los Asentamientos Humanos (ONU-Hábitat). Según los
metadatos oficiales, actualizados el 23 de abril de 2025, el acceso se considera conveniente
cuando una parada oficialmente reconocida es accesible a una distancia a pie de 500~m,
medida a lo largo de la red vial, desde un punto de referencia como el hogar, la escuela, el
lugar de trabajo o el mercado, en el caso de sistemas de baja capacidad (por ejemplo,
autobús o BRT), o de 1~km en el caso de sistemas de alta capacidad (tren, metro,
transbordador) \parencite{onu2025}. Los metadatos añaden criterios complementarios, como la
accesibilidad para personas con necesidades especiales, la frecuencia del servicio en horas
punta y la seguridad y comodidad del entorno de la parada.

Este umbral se ha operacionalizado ampliamente con datos abiertos. El Center for
International Earth Science Information Network (CIESIN, \citeyear{ciesin2023}) calculó el
indicador para 5749 centros urbanos de 178 países, como proporción de la población de
WorldPop situada a no más de 0,5~km a pie de una parada de baja capacidad registrada en
OpenStreetMap, y oficinas estadísticas nacionales como \textcite{statcan2023} han aplicado
el mismo estándar de 500~m a lo largo de la red vial. En la literatura técnica
norteamericana, el \textit{Transit Capacity and Quality of Service Manual} (TCQSM), en su
tercera edición \parencite{tcqsm2013}, ofrece un marco para medir la disponibilidad del
transporte público desde la perspectiva del pasajero, incluida la cobertura espacial del
servicio, en el que se emplean distancias de caminata del orden de 400~m (un cuarto de
milla) para paradas de autobús y de hasta 800~m (media milla) para estaciones de modos de
alta capacidad. El rango de 400 a 800~m constituye así un intervalo de referencia razonable
para la distancia que los usuarios están dispuestos a caminar hasta el transporte público.

La adaptación de este concepto al paratránsito exige una decisión explícita. Dado que en el
sistema cochabambino el ascenso puede producirse en cualquier punto del recorrido, el
conjunto de puntos de acceso potenciales no es un conjunto de paradas, sino la propia línea
del trazado. Por ello, en el proyecto se mide la cobertura como la distancia de la zona al
trazado más cercano de las rutas del feed (\var{shapes.txt}) y se adopta el umbral de 500~m
del indicador 11.2.1 como criterio operativo de \enquote{zona cubierta}. Esta adaptación
introduce dos simplificaciones que deben declararse: se sustituye la parada oficialmente
reconocida por el trazado, y la distancia por red vial por una distancia geométrica, lo que
tiende a sobrestimar la cobertura allí donde la red vial es discontinua. El umbral es,
además, de una magnitud comparable a la arista media de una celda H3 de resolución~8
($\approx 0{,}53$~km; véase \ref{subsec:h3}), lo que hace coherente la escala de la
medida con la unidad de análisis. La cobertura así definida es un indicador de oferta
mapeada, no de accesibilidad efectiva, y su sensibilidad al umbral debe explorarse.

\subsection{Datos de consultas de una aplicación de ruteo como señal de demanda de información}
\label{subsec:consultas-senal}

La variable de interés del proyecto es el número de consultas de ruta registradas por ``Trufi App'', cada una caracterizada por un origen, un destino, una marca temporal y un identificador
de instalación. Conviene subrayar que una consulta no es un viaje. Una consulta expresa una
necesidad de información sobre cómo desplazarse entre dos puntos: un mismo viaje puede
generar varias consultas (por ejemplo, al comparar alternativas o repetir la búsqueda), y
muchos viajes no generan ninguna, en particular los desplazamientos habituales cuyos
usuarios ya conocen la ruta. Por ello, las consultas constituyen una señal de demanda de
información de movilidad y, solo de forma indirecta y sesgada, una señal de demanda de
viajes.

Esta señal está sujeta a varios sesgos. El primero es el de adopción: solo generan
consultas quienes tienen un teléfono inteligente, han instalado la aplicación y la usan, y
la adopción varía según el nivel socioeconómico, la edad y la exposición a campañas de
difusión; la propia Trufi Association atribuye parte del crecimiento de instalaciones al
lanzamiento de una nueva versión y a la recomendación espontánea de la creadora de
contenido de TikTok Carla Salló, con más de 1,5~millones de seguidores, tras la cual la
asociación registró un fuerte aumento de las solicitudes de itinerarios desde la aplicación
de Cochabamba \parencite{trufi-20000}. El segundo es el de necesidad de información: los
viajes poco familiares (hacia destinos nuevos, por visitantes o por trámites ocasionales)
generan proporcionalmente más consultas que los rutinarios. El tercero es el de
disponibilidad de oferta en la aplicación: si una zona no está servida por rutas mapeadas,
la aplicación no ofrece itinerarios útiles, lo que puede reducir su uso y crear un círculo
en el que la falta de cartografía y la escasez de consultas se refuerzan mutuamente.

La literatura sobre datos digitales pasivos y datos de aplicaciones como indicadores de
movilidad ha documentado ampliamente estos problemas de representatividad.
\textcite{wesolowski2013} cuantificaron, a partir de los movimientos diarios de casi
15~millones de personas en Kenia a lo largo de un año, el efecto del sesgo de tenencia de
teléfonos en las estimaciones de movilidad, y mostraron que las inferencias a nivel
poblacional se complican porque la tenencia difiere entre grupos demográficos con movilidad
distinta. \textcite{blondel2015}, en su revisión de resultados con datos de telefonía
móvil, subrayan tanto el potencial de estos datos para el análisis urbano como la necesidad
de tratar sus sesgos de cobertura. En consecuencia, las consultas de ``Trufi App'' no pueden
interpretarse como medida absoluta de la demanda de transporte, pero sí como señal relativa
cuyo nivel esperado puede modelarse en función de la población y del contexto territorial.

La implicación metodológica es que el estimando del proyecto se define explícitamente como
el número esperado de consultas, y no de viajes, de cada zona dada su población y su
contexto. En este marco, una zona poblada con cero consultas no es un dato faltante que deba
descartarse, sino una observación informativa, porque su contraste con lo esperado es
precisamente lo que puede señalar un déficit de cartografía o de adopción. Esta decisión
diferencia al proyecto del antecedente de \textcite{angulo2024}, monografía de la UMSS que
abordó la predicción de demanda en ``Trufi App'' con la metodología CRISP-DM y que excluyó las
zonas con datos insuficientes.

\subsection{Población como exposición: Kontur Population y fuentes alternativas}
\label{subsec:kontur}

La población de cada zona se toma de Kontur Population, un conjunto de datos mundial de
población distribuida en hexágonos vectoriales H3 de resolución~8 (descritos por su
productor como hexágonos de \enquote{400~m}), publicado en la plataforma Humanitarian Data
Exchange (HDX) bajo licencia Creative Commons Atribución Internacional
\parencite{kontur2023}. La ficha de HDX lo describe como una fusión depurada de datos de
GHSL, Facebook, Microsoft Buildings, Copernicus Global Land Service Land Cover, Land
Information New Zealand y OpenStreetMap, y precisa que el cálculo de la población se basa en
superponer la capa Global Human Settlement Layer (GHSL) con la población de la High
Resolution Settlement Layer (HRSL) de Facebook allí donde esta está disponible,
restringiendo con datos de OpenStreetMap los artefactos conocidos de ambas fuentes. La
documentación del productor añade que las huellas de edificios de Microsoft, los datos de
Land Information New Zealand y la cobertura del suelo de Copernicus se emplean para mejorar
la precisión de la distribución; que las canteras y las vías de gran tamaño se marcan como
no pobladas, porque GHSL las detecta a menudo erróneamente como pobladas, al igual que
lagos, ríos, glaciares, arenales y bosques; que las celdas extremadamente pobladas se
redistribuyen hacia celdas vecinas para satisfacer restricciones, preservando el total; y
que los conteos no enteros se redondean \parencite{kontur-sf}. En HDX figuran ediciones
fechadas entre el 11 de marzo de 2020 y el 1 de noviembre de 2023. El proyecto emplea la
edición 2023-11-01, la más reciente; la de 2022-06-30 no se incorporó a los datos de este
trabajo y queda como extensión para el análisis de sensibilidad
(sección~\ref{sec:interpretacion-resultados}).

Kontur Population es, por tanto, una estimación modelada de tipo dasimétrico, y no un conteo
censal. Su lógica consiste en redistribuir totales de población de fuentes de base sobre la
superficie a partir de indicios de ocupación (superficie construida, edificios, cobertura
del suelo), lo que produce buenas aproximaciones agregadas pero errores locales no
despreciables en celdas individuales, sobre todo en áreas periurbanas de crecimiento
reciente o de edificación dispersa. La referencia oficial de población en Bolivia es el
Censo de Población y Vivienda 2024 del Instituto Nacional de Estadística (INE), realizado el
23 de marzo de 2024, que registró una población nacional de 11\,365\,333 habitantes
\parencite{ine2025}; en la socialización de los resultados del conteo poblacional, el INE
informó que el departamento de Cochabamba registró 2\,005\,373 habitantes
\parencite{ine2024}. La edición de Kontur empleada es anterior a la publicación de
estos resultados censales y se apoya en insumos de fechas previas, por lo que puede no
reflejar plenamente la distribución de 2024; esta discrepancia temporal justifica tratar la
población como una covariable medida con error y evaluar la sensibilidad de las estimaciones
a la edición empleada.

Existen otras fuentes de población en grilla que conviene conocer para situar la elección.
WorldPop produce estimaciones de alta resolución mediante modelos de desagregación de datos
censales con bosques aleatorios alimentados por covariables geoespaciales
\parencite{stevens2015,tatem2017}, y GHS-POP, del Centro Común de Investigación de la
Comisión Europea, distribuye la población sobre la superficie construida detectada por
satélite \parencite{schiavina2023}. Las diferencias metodológicas entre estas fuentes se
traducen en discrepancias locales apreciables, por lo que ninguna puede considerarse verdad
de terreno a escala de barrio. Kontur presenta, para este proyecto, la ventaja práctica de
estar publicado de forma nativa sobre la grilla H3 de resolución~8, lo que evita una
reproyección y una reagregación adicionales.

Por último, debe distinguirse entre población residente y población flotante. Kontur estima
la población residente, es decir, dónde vive la gente, mientras que las consultas pueden
originarse también en lugares de actividad diurna (centros comerciales, mercados,
universidades, zonas de oficinas o terminales) cuya población residente es baja. En
consecuencia, en las zonas centrales o de alta actividad cabe esperar más consultas de las
que su población residente predeciría, y en zonas dormitorio la relación puede invertirse.
La implicación es que la población actúa como medida de exposición necesaria pero no
suficiente, y que las variables de contexto territorial (distancia al centro, municipio,
vecindad) cumplen la función de absorber parte de la diferencia entre población residente y
actividad efectiva.

\subsection{La lógica \texorpdfstring{\enquote{observado frente a esperado}}{"observado frente a esperado"} como instrumento de priorización}
\label{subsec:observado-esperado}

La comparación entre un conteo observado y un conteo esperado es un procedimiento clásico
de la epidemiología. La razón estandarizada de mortalidad o de morbilidad (SMR, por sus
siglas en inglés) de un área~$i$ se define como $\mathrm{SMR}_i = O_i / E_i$, donde $O_i$ es
el número de casos observados y $E_i$ el número esperado si el área tuviera las tasas de una
población de referencia; en la estandarización indirecta, $E_i = \sum_k P_{ik}\,\lambda_k$,
donde $P_{ik}$ es la población del estrato~$k$ en el área~$i$ y $\lambda_k$ la tasa de
referencia del estrato. Una razón superior a uno indica más casos de los esperados, y una
inferior, menos. La lógica es transferible a cualquier conteo cuya magnitud depende de una
población en riesgo: no interesa el conteo bruto, que refleja sobre todo el tamaño de la
población, sino su desviación respecto de lo que la población y la estructura de referencia
harían esperar.

La cartografía de enfermedades puso de relieve una dificultad central de este enfoque, el
problema de los números pequeños. \textcite{clayton1987} mostraron que, en áreas con
poblaciones pequeñas, las razones $O/E$ son muy inestables: un solo caso adicional puede
duplicar la razón, de modo que los valores extremos del mapa tienden a concentrarse en las
áreas menos pobladas por puro efecto de la varianza muestral. Su propuesta consistió en
estimadores empírico-bayesianos que contraen cada razón hacia un valor de referencia en
proporción inversa a la información disponible, y \textcite{marshall1991} desarrolló
estimadores empírico-bayesianos, globales y locales, para cartografiar tasas de enfermedad y
mortalidad con ese mismo propósito. Este problema es directamente pertinente para el
proyecto, porque un conjunto de zonas H3 de resolución~8 incluye inevitablemente zonas de
población reducida cuyas razones $O/E$ serían engañosamente extremas.

La transposición al problema de ``Trufi App'' es natural: se sustituye el caso de enfermedad
por la consulta de ruta, la población en riesgo por la población residente de la zona y la
tasa de referencia por el modelo que estima las consultas esperadas dado el contexto. La
brecha de cada zona puede expresarse en términos absolutos, $B_i = O_i - E_i$, en términos
relativos, $R_i = O_i / E_i$, o en términos probabilísticos, como la probabilidad de
observar $O_i$ o menos consultas si el verdadero valor esperado fuera $E_i$,
$p_i = P(Y \le O_i \mid \mu = E_i)$. Esta última forma incorpora la incertidumbre muestral y
mitiga el problema de los números pequeños, porque una misma razón es más significativa
cuando el valor esperado es grande. Una zona con muchas consultas esperadas y pocas
observadas se convierte así en candidata prioritaria para la revisión cartográfica.

Es imprescindible, no obstante, subrayar que la brecha orienta la revisión pero no
diagnostica por sí sola la falta de rutas. Un déficit de consultas respecto de lo esperado
puede deberse a rutas existentes pero no mapeadas, a rutas mapeadas con errores, a una baja
adopción de la aplicación en la zona, a errores en la estimación de población, a usos del
suelo no residenciales o a limitaciones del propio modelo. La brecha actúa, por tanto, como
un filtro que ordena el territorio según la magnitud de la anomalía y reduce el espacio de
búsqueda de los voluntarios, mientras que la confirmación de la causa corresponde a la
verificación cartográfica en OSM y en terreno.
